% ============================================================
% MÓDULO 6: TEMPORAL-DIFFERENCE LEARNING (SARSA Y Q-LEARNING)
% ============================================================
\section{M6 · TD Learning}
\module{Módulo 6: TD Learning}{La combinación definitiva: muestreo sin modelo + estimación por bootstrapping}

\begin{frame}{La Idea Central de la Diferencia Temporal: Bootstrapping}
  \begin{columns}[c,onlytextwidth]
    \column{0.50\textwidth}
      \textbf{¿Qué es Bootstrapping?}
      Actualizar una predicción a partir de \textbf{otra predicción futura}, sin esperar al desenlace final del episodio.
      \vspace{0.25cm}
      \begin{itemize}
        \item \textbf{DP:} Hace bootstrapping, pero necesita el modelo $\mathcal{P}$ y $\mathcal{R}$.
        \item \textbf{MC:} No hace bootstrapping (espera al final), pero no necesita modelo.
        \item \textbf{TD:} \textbf{¡Lo mejor de ambos mundos!} Aprende de la experiencia pura y actualiza paso a paso.
      \end{itemize}
    \column{0.48\textwidth}
      \begin{beamercolorbox}[rounded=true,sep=0.8em]{block body}
        \centering
        {\large $V(S_t) \leftarrow V(S_t) + \alpha \big[ \underbrace{R_{t+1} + \gamma V(S_{t+1})}_{\text{Objetivo TD}} - V(S_t) \big]$}
      \end{beamercolorbox}
      \vspace{0.25cm}
      \small El \textbf{Error TD} $\delta_t = R_{t+1} + \gamma V(S_{t+1}) - V(S_t)$ mide la sorpresa entre lo que predije ayer y la nueva evidencia de hoy.
  \end{columns}
  \vspace{0.35cm}
  \keyidea{TD actualiza en cada paso individual ($t \to t+1$): puede aprender en tareas continuas sin fin de turno.}
\end{frame}

\begin{frame}{Control: SARSA vs. Q-Learning}
  \begin{columns}[T,onlytextwidth]
    \column{0.48\textwidth}
      \begin{block}{SARSA (On-Policy)}
        Aprende el valor de la política que realmente ejecuta (incluyendo los errores de exploración):
        \[
          Q(S,A) \leftarrow Q(S,A) + \alpha [ R + \gamma Q(S', A') - Q(S,A) ]
        \]
        \begin{itemize}
          \item La tupla: $(S_t, A_t, R_{t+1}, S_{t+1}, A_{t+1})$.
          \item $A'$ se elige con la misma política $\varepsilon$-greedy.
          \item Tiende a ser más \textbf{prudente y conservador}.
        \end{itemize}
      \end{block}
    \column{0.48\textwidth}
      \begin{alertblock}{Q-Learning (Off-Policy, Watkins 1989)}
        Aprende la política óptima $Q^*$ sin importar qué política de exploración use el robot:
        \[
          Q(S,A) \leftarrow Q(S,A) + \alpha [ R + \gamma \max_{a'} Q(S', a') - Q(S,A) ]
        \]
        \begin{itemize}
          \item El objetivo usa $\max_{a'}$, no la acción ejecutada.
          \item Aprende la ruta \textbf{óptima más corta}.
        \end{itemize}
      \end{alertblock}
  \end{columns}
\end{frame}

\begin{frame}{Cálculo Manual Paso a Paso: Actualización de Q-Learning}
  \begin{columns}[c,onlytextwidth]
    \column{0.50\textwidth}
      \textbf{Transición observada de OptiBot:}
      \begin{itemize}
        \item Estado actual: $S = \text{Pasillo A}$
        \item Acción: $A = \text{Avanzar}$
        \item Recompensa recibida: $R = -1.0$
        \item Siguiente estado: $S' = \text{Pasillo B}$
      \end{itemize}
      \textbf{Valores previos en la tabla Q:}
      \begin{itemize}
        \item $Q(S, A) = 4.0$
        \item En $S'$: $Q(S', \text{Izq}) = 6.0$, $Q(S', \text{Der}) = 8.0$
        \item Por tanto: $\max_{a'} Q(S', a') = 8.0$
        \item Parámetros: $\alpha = 0.20$, $\gamma = 0.90$
      \end{itemize}
    \column{0.48\textwidth}
      \begin{exampleblock}{Cálculo en la pizarra}
        \begin{align*}
          \text{Objetivo TD} &= R + \gamma \max_{a'} Q(S', a') \\
          &= -1.0 + 0.90(8.0) = \mathbf{6.20} \\
          \delta_t &= \text{Objetivo} - Q(S, A) \\
          &= 6.20 - 4.00 = \mathbf{+2.20} \\
          Q_{\text{nuevo}}(S, A) &= 4.00 + 0.20(+2.20) \\
          &= 4.00 + 0.44 = \mathbf{4.44}
        \end{align*}
      \end{exampleblock}
      \small ¡La estimación subió de 4.00 a 4.44 al descubrir que el Pasillo B tiene un camino muy prometedor!
  \end{columns}
\end{frame}

\begin{frame}[fragile]{Implementación de Q-Learning para el Almacén}
  \begin{columns}[T,onlytextwidth]
    \column{0.62\textwidth}
\begin{lstlisting}[style=rlcode]
for ep in range(num_episodes):
    s, _ = env.reset()
    done = False
    
    while not done:
        # Seleccion de accion epsilon-greedy
        if np.random.rand() < epsilon:
            a = env.action_space.sample()
        else:
            a = np.argmax(Q[s])
            
        s_next, r, terminated, truncated, _ = env.step(a)
        done = terminated or truncated
        
        # Actualizacion de Q-Learning
        target = r
        if not done:
            target += gamma * np.max(Q[s_next])
            
        td_error = target - Q[s, a]
        Q[s, a] += alpha * td_error
        s = s_next
\end{lstlisting}
    \column{0.35\textwidth}
      \begin{alertblock}{La limitación de M6}
        Q-Learning tabular es el estándar de oro para problemas pequeños.\\
        \vspace{0.2cm}
        \textbf{El abismo de la dimensionalidad:}
        Si el almacén tiene $1,000,000$ posiciones o sensores continuos (LiDAR, batería en \%, velocidades), la tabla requeriría gigabytes y nunca visitaría cada casilla.\\
        \vspace{0.2cm}
        ¿Cómo generalizar a estados nunca vistos? $\to$ \textbf{Deep Q-Networks (M7)}.
      \end{alertblock}
  \end{columns}
\end{frame}
